% -----------------------------------------------------------------------------------
\section{Optimality criteria and the possibility of ex-ante optimality}\label{sec:criteria}

\subsection{Within-problem criteria}

The criteria below are the precise forms of the principles that the main decision theories are
usually taken to satisfy. Each is a condition on the policy $T_D$ that a theory $T$ selects in a
problem $D$; the commitment criterion \cri{4} and the evidential criterion \cri{6} are introduced in
Section~\ref{sec:commit}, once the constructions they require are available.

\begin{definition}[Ex-ante optimality, \cri{1}]
$T$ satisfies \cri{1} on $D$ if \[ T_D\in\argmax_{\pi\in\Pol_D}V_D(\pi). \]
\end{definition}

\begin{definition}[Structural causal dominance, \cri{2}]
$T$ satisfies \cri{2} on $D$ if, at every information set $h$ positively reached under $\pi=T_D$,
there is no $a\in A_h$ such that
\[
  U_D\bigl(s,\dev{\pi}{h}{a}\bigr)>U_D\bigl(s,\dev{\pi}{h}{\pi(h)}\bigr)
  \quad\text{for every }s\in R_h(\pi).
\]
\end{definition}

The quantifier ranges over \emph{structural} states: all plans that reach $h$, including those to
which $\mu_\pi$ assigns probability zero. This is a deliberate resolution of an ambiguity in
\enquote{every state compatible with $h$}, and it matters. If compatibility meant only the states in
the support of $\mu_\pi$, then in Death in Damascus with an infallible predictor (Section~\ref{sec:did})
the single supported state places Death in the chosen city, so switching would strictly improve the
payoff in every supported state and \cri{2} itself would be unsatisfiable. Likewise, on trees where
prediction nodes follow $h$, \enquote{hold fixed the outcomes of the nodes preceding $h$} leaves open
how the later prediction nodes are to be evaluated; freezing their policy dependence while
averaging over their outcomes can make dominance collapse into an unsatisfiable best-response
requirement. The structural plan fixes both matters at once.

\begin{definition}[Ratifiability, \cri{3}]
$T$ satisfies \cri{3} on $D$ if, at every information set $h$ positively reached under $\pi=T_D$,
\[
  \pi(h)\in\argmax_{a\in A_h}Q_\pi(h,a).
\]
\end{definition}

\begin{definition}[Existence, \cri{5}]
$T$ satisfies \cri{5} on $D$ if its own criterion is satisfiable there, that is, if $T_D$ is defined.
\end{definition}

Since decision theories are total by Definition~\ref{def:theory}, \cri{5} holds for every theory
and imposes no restriction; it records the requirement that a criterion be satisfiable on every
problem of the class. Were partial theories admitted, totality would become an independent axiom
and the consistency maps of Section~\ref{sec:commit} would change.

\begin{proposition}\label{prop:c3c2}
\cri{3} implies \cri{2}.
\end{proposition}

\begin{proof}
Suppose $h$ is positively reached under $\pi=T_D$ and some $a$ satisfies
$U_D(s,\dev{\pi}{h}{a})>U_D(s,\pi)$ for every $s\in R_h(\pi)$. Since $r_\pi(h)>0$, the conditional
distribution $\mu_\pi(\cdot\mid R_h(\pi))$ is a probability measure carried by $R_h(\pi)$, and a strict
pointwise inequality on $R_h(\pi)$ gives $Q_\pi(h,a)>Q_\pi(h,\pi(h))$. Thus $\pi(h)$ is not a maximiser
of $Q_\pi(h,\cdot)$, and \cri{3} fails.
\end{proof}

\subsection{Ex-ante optimality is always attainable}

The opening question admits a direct answer once \enquote{decision theory} is taken in the
permissive sense of Definition~\ref{def:theory}.

\begin{theorem}[Universal ex-ante selection]\label{thm:universal}
On every class of well-defined finite problems there is a decision theory satisfying \cri{1} on
every problem of the class.
\end{theorem}

\begin{proof}
Fix a tie-breaking order on each $\Pol_D$ and set
$T^{\mathrm{ex}}_D=\first\argmax_{\pi\in\Pol_D}V_D(\pi)$. Each $\Pol_D$ is finite and nonempty, and
$V_D$ is real valued on it because $D$ is well defined, so the maximum exists.
\end{proof}

\begin{remark}[Self-referential problems]\label{rem:skalse}
\textcite{skalse2021} proposes, for a fixed target theory $X$, a choice between a transparent box
containing $1000$ and an opaque box whose content is $0$ if $X$ permits choosing it in that problem and
$2000$ otherwise, and argues that another theory can outperform $X$. The target theory is an
environmental constant, so this is not discrimination by the identity of the evaluated agent.
Translated into the present model, however, the construction requires a fixed point: for an ex-ante
optimiser, an empty opaque box makes it choose the transparent box, which requires the opaque box to
be full, and conversely. Neither content yields a well-defined instance. In the same way, the
diagonal family $D_b$ with $u(a)=\one\{a\neq b\}$, $b\in\{0,1\}$, would need $b=T(D_b)$, but an ex-ante
optimiser has $T(D_0)=1$ and $T(D_1)=0$. The schema does not refute Theorem~\ref{thm:universal} on
well-defined problems; its additional self-referential realisability assumption is where it leaves
that class. \textcite{lewis1981rich} analysed a related failure: an announced reward for supposedly
irrational choice can render the proposed scenario inconsistent. The verdict depends on the problem
being a finite object whose data include the constant's output. When problems are given as programs,
so that the agent must compute that output itself, the same construction is well defined and defeats
every agent that halts on all inputs (Theorem~\ref{thm:fixedcopy}).
\end{remark}

The unqualified thesis that no decision theory is optimal on all fair problems is therefore false.
What can be true is that particular \emph{combinations} of criteria are jointly unsatisfiable, and
the remaining sections determine exactly which. A first observation shows how little such a result
can depend on generality or uniformity. Call a criterion $C$ \emph{local} if whether a theory $T$
satisfies $C$ on $D$ depends only on $T_D$; \cri{1}, \cri{2}, \cri{3}, and \cri{5} are local, while
\cri{4} is not. For a local criterion $C_i$ let $A_i(D)\subseteq\Pol_D$ be the set of policies that
satisfy it on $D$.

\begin{lemma}[Pointwise reduction]\label{lem:pointwise}
For local criteria $C_1,\dots,C_k$ and a class $\cls$,
\[
  \exists T\;\forall D\in\cls\;\bigwedge_{i}C_i(T,D)
  \quad\Longleftrightarrow\quad
  \forall D\in\cls\;\;\bigcap_{i}A_i(D)\neq\varnothing.
\]
\end{lemma}

\begin{proof}
From left to right, $T_D$ lies in the intersection. From right to left, let $T_D$ be the first
element of the intersection under a fixed order.
\end{proof}

Because lookup tables are admitted, an impossibility for local criteria is never more than the
exhibition of a single problem on which their acceptable sets are disjoint. No difficulty about
general rules, computability, or self-reference is involved; those enter only when a theory must be
a single rule on a common description of problems (Part~\ref{part:across}) or a computable procedure
that reads descriptions (Part~\ref{part:programs}). This is the precise sense in which the
within-problem impossibilities of the next two sections are shallow; non-local criteria, and the
cross-problem axioms of Sections~\ref{sec:matrix}--\ref{sec:exante}, are where deeper structure
can appear.
